\documentclass[11pt,a4paper]{article}
\usepackage{xepersian}
\settextfont{Arial}\setlatintextfont{Times New Roman}
\author{مجید فیض دیلمی}\date{}
\begin{document}
\# ArchGuard-APR: تعمیر خودکار برنامه آگاه از قیود معماری با استفاده از گراف‌های مخزن و بازخورد اجرا

**مجید فیض دیلمی**<br>
دانشجوی کارشناسی ارشد مهندسی نرم‌افزار<br>
گروه مهندسی کامپیوتر، واحد لاهیجان، دانشگاه آزاد اسلامی، لاهیجان، ایران<br>
مسئول مکاتبات: majid.feiz@iau.ir · ORCID: https://orcid.org/0009-0006-8439-1839<br>
*نسخه نویسنده — داوری علمی نشده*

\#\# چکیده

تعمیر خودکار برنامه (Automated Program Repair) در مقیاس مخزن باید افزون بر رفتار موضعی، ساختار
سامانه را نیز حفظ کند. یک تغییر پیشنهادی ممکن است آزمونی محدود را بگذراند، اما مرز وابستگی را نقض
کند، مسیر ناایمنی را تغییر دهد، یا بر زمینه‌ای ناقص تکیه کند. ArchGuard-APR نمونه‌ای پژوهشی برای
تعمیر آگاه از معماری است که تحلیل قطعی مخزن، گراف‌های تغییرناپذیر، پروفایل‌های صریح معماری، بازیابی
زمینه مبتنی بر گراف، هماهنگی نقش‌محور، کنترل سیاست وصله، بازخورد اجرای ایزوله، محدودیت هزینه و شواهد
استاندارد و عاری از کد منبع را یکپارچه می‌کند. ارزیابی انجام‌شده یک مطالعه زوجی اکتشافی امکان‌پذیری
است و نه آزمایشی تأییدی درباره اثربخشی. از شش اجرای برنامه‌ریزی‌شده، چهار اجرا کامل شد و دو وضعیت
black-4 پیش از تماس با ارائه‌دهنده دچار ریزش زیرساختی شدند. در دو وظیفه زوجی معتبر، نرخ موفقیت خط
پایه و ArchGuard هر دو ۰ از ۲ بود و هر چهار خروجی کامل‌شده «وصله نامعتبر» طبقه‌بندی شدند. چهار اجرا
۳۵٬۴۳۱ توکن ورودی و ۲٬۷۳۷ توکن خروجی، در مجموع ۳۸٬۱۶۸ توکن، با هزینه ۰٫۱۲۹۶۳۲۵ دلار مصرف کردند.
تفاوتی در اثربخشی مشاهده نشد و اندازه نمونه برای استنباط آماری کافی نیست. نتیجه اصلی، امکان‌پذیری
مهندسی خط لوله و ضرورت بهبود اعتبار وصله و پایداری زیرساخت است.

**واژگان کلیدی—** تعمیر خودکار برنامه؛ معماری نرم‌افزار؛ گراف مخزن؛ قید معماری؛ مدل زبانی بزرگ؛
سامانه چندعاملی؛ بازخورد اجرا؛ بازتولیدپذیری.

\#\# ۱. مقدمه

تعمیر خودکار برنامه می‌کوشد تغییراتی تولید کند که نقص نرم‌افزار را برطرف سازند. در وظایف سطح مخزن،
مسئله به یک تابع محدود نیست: ابزار باید کد مرتبط را بیابد، سازمان پروژه را رعایت کند، وصله‌ای مجاز
بسازد و رفتار را در محیطی کنترل‌شده اعتبارسنجی کند. عامل‌های مبتنی بر مدل زبانی امکان برنامه‌ریزی و
کار با ابزارها را افزایش می‌دهند، اما مرزهای شکست بیشتری نیز پدید می‌آورند. پاسخ متنی محتمل لزوماً
وصله‌ای قابل تجزیه نیست؛ وصله قابل اعمال نیز الزاماً ایمن یا منطبق با معماری نیست.

ArchGuard-APR این مرزها را اشیای صریح پژوهش می‌داند. هسته دامنه و کاربرد از سازگارسازهای زیرساختی و
فروشنده مستقل است. منبع مخزن ابتدا به تصویر لحظه‌ای اعتبارسنجی‌شده تبدیل می‌شود؛ ساخت گراف تنها همان
تصویر را مصرف می‌کند؛ ارزیابی معماری نیز فقط گراف معتبر و پروفایل نسخه‌دار را می‌پذیرد. بازیابی زمینه
بودجه‌دار است و بازرسی وصله و اجرا زیر سیاست‌های ایمنی، ایزوله‌سازی و ثبت قطعی شواهد انجام می‌شود.

مقاله دو سهم جدا دارد. سهم مهندسی، نمونه‌ای پیاده‌سازی‌شده از تحلیل مخزن، قیود معماری، بازیابی، نقش‌های
تعمیر، اعتبارسنجی وصله و بازخورد اجراست. شواهد علمی محدودترند: چهار خروجی نامعتبر و دو ریزش زیرساختی.
بنابراین مقاله ادعای برتری ندارد، آزمون معنادار استنباطی انجام نمی‌دهد و نتایج منفی را حذف نمی‌کند.

\#\# ۲. مبانی و انگیزه

روش‌های تولید و اعتبارسنجی، نامزدها را می‌سازند و در برابر نشانه‌ای از درستی می‌آزمایند. در مقیاس
مخزن، زمینه مرتبط ممکن است از مرز ماژول عبور کند، محیط ساخت مستقل از وصله شکست بخورد و مجموعه آزمون
مشخصاتی ناقص باشد. معماری بعد دیگری از درستی را معرفی می‌کند: مرزهای وابستگی و مالکیت می‌توانند قصد
نگهداشت‌پذیری را بیان کنند که آزمون رفتاری آن را نمی‌سنجد.

\#\#\# ۲.۱. خط لوله تعمیر و معنای پیامدها

GenProg جست‌وجوی مبتنی بر آزمون را به کار می‌گیرد [4]؛ SemFix از تحلیل معنایی [5] و Angelix از سنتز
نمادین چندخطی استفاده می‌کند [7]. Prophet رتبه‌بندی نامزدها را از وصله‌های انسانی می‌آموزد [8] و
Elixir زمینه شیءگرا را وارد انتخاب می‌کند [10]. این تفاوت‌ها نشان می‌دهند سازوکار پیشنهاد وصله از
داوری پذیرش آن جداست. مطالعات Defects4J نیز نشان داده‌اند وصله آزمون‌گذر می‌تواند نادرست بماند [9].

در این مقاله «نامزد» هر ویرایش پیشنهادی است؛ «وصله معتبر» قرارداد قالب، نحو، مسیر، اندازه و سیاست را
رعایت می‌کند؛ و موفقیت تعمیر افزون بر قابلیت اعمال، به نتیجه ارزیابی وابسته است. پاسخ ارائه‌دهنده که
به وصله مجاز تبدیل نشود یک اجرای علمی با طبقه‌بندی `invalid_patch` است. شکست کسب یا اعتبارسنجی مخزن
پیش از ارسال به ارائه‌دهنده، پیامد مدل نیست و ریزش زیرساختی ثبت می‌شود.

\#\#\# ۲.۲. تعمیر سطح مخزن و زمینه

Defects4J مجموعه‌ای کنترل‌شده از خطاهای جاوا [6] و BugsInPy مجموعه‌ای متناظر برای پایتون فراهم می‌کند
[14]. SWE-bench مسئله را به رفع درخواست واقعی در مخزن و محیط اجرای آن گسترش می‌دهد [20]. این
بنچمارک‌ها از نظر زبان، وابستگی، ساخت وظیفه و محیط یکسان نیستند. بازیابی افزوده، شواهد بازیابی‌شده را
از پارامترهای مدل جدا می‌کند [15]؛ نمایش‌های گرافی برنامه [11] و GraphCodeBERT [16] نیز سودمندی
روابط صریح را نشان می‌دهند. گراف ایستای این سامانه با نمایش‌های یادگرفته‌شده یکسان نیست؛ هدف آن انتخاب
محدود، قابل ممیزی و تکرارپذیر زمینه است.

\#\#\# ۲.۳. معماری و انطباق

معماری نرم‌افزار با عناصر، روابط و منطق طراحی توصیف می‌شود [1]. مدل‌های بازتابی فاصله میان مدل سطح
بالا و ساختار منبع را می‌سنجند [2]. در ArchGuard-APR پروفایل بازبینی‌شده انسانی، انتساب ماژول‌ها و
وابستگی‌های مجاز را تعیین می‌کند. سامانه معماری را استنباط نمی‌کند و در برابر ماژول مبهم یا بدون
انتساب، بسته و محافظه‌کارانه شکست می‌خورد.

\#\# ۳. کارهای مرتبط

مرور Monperrus گستره تعمیر رفتاری و انواع داور را صورت‌بندی می‌کند [3]. CoCoNuT ترجمه عصبی آگاه از
زمینه را به کار می‌گیرد [13] و ارزیابی مدل‌های بزرگ از ظرفیت آن‌ها برای تولید وصله گزارش می‌دهد [19].
RepairThemAll تفاوت ابزار و بنچمارک را در مقیاس بزرگ مستند کرده است [12]. این مطالعات اهمیت فضای
نامزد، زمینه و اعتبارسنجی را روشن می‌کنند، اما اثربخشی ArchGuard-APR را ثابت نمی‌کنند.

در تعمیر مخزن، RepairAgent حلقه ابزارمحور [21]، AutoCodeRover جست‌وجوی ساختاریافته [22] و SpecRover
استخراج قصد و نقش بازبین [23] را بررسی می‌کنند. ReAct استدلال و عمل را درهم می‌آمیزد [17] و Reflexion
از بازخورد زبانی استفاده می‌کند [18]. سامانه حاضر از این ایده‌ها برای جداسازی نقش‌ها الهام می‌گیرد،
ولی نوآوری یا تقدم تاریخی فراگیر ادعا نمی‌کند. تمایز آن در ترکیب پروفایل معماری صریح با شواهد قطعی،
کنترل ایمنی و حسابداری اجرای منبع‌آزاد است.

\#\#\# ۳.۱. تعمیر تولید و اعتبارسنجی، معنایی و مبتنی بر یادگیری

خانواده تولید و اعتبارسنجی، فضای تغییر را جست‌وجو و هر نامزد را با یک داور اجرایی می‌سنجد. GenProg
نمونه شناخته‌شده این الگوست [4]، درحالی‌که SemFix و Angelix محدودیت‌های معنایی و سنتز نمادین را برای
کاهش فضای جست‌وجو به کار می‌گیرند [5]، [7]. Prophet از وصله‌های انسانی برای رتبه‌بندی می‌آموزد [8]
و Elixir شواهدی مانند زمینه شیءگرا را وارد انتخاب می‌کند [10]. این مسیر پژوهشی نشان می‌دهد تولید،
رتبه‌بندی و داوری سه مسئولیت متمایزند؛ موفقیت در یکی جایگزین اعتبارسنجی دیگری نیست.

روش‌های عصبی مانند CoCoNuT زمینه را در تولید یادگرفته‌شده وارد می‌کنند [13] و مطالعه مدل‌های زبانی
بزرگ ظرفیت و محدودیت آن‌ها را برای تولید وصله بررسی می‌کند [19]. ArchGuard-APR رقیب مستقیم هیچ‌یک
از این مولدها نیست. ارائه‌دهنده پشت یک درگاه قابل تعویض قرار دارد و سهم سامانه، قرارداد پیرامونی برای
زمینه، معماری، سیاست وصله، اجرا و شواهد است. ازاین‌رو نتایج فاز ۱۳ درباره برتری الگوریتم تولید یا مدل
خاصی نیست.

\#\#\# ۳.۲. بنچمارک‌ها و حل مسئله در سطح مخزن

Defects4J خطاهای واقعی جاوا را همراه زیرساخت کنترل‌شده فراهم می‌کند [6] و BugsInPy همین هدف را برای
پروژه‌های پایتون دنبال می‌کند [14]. RepairThemAll نشان داده است که اجرای یکسان ابزارها بر بنچمارک‌های
متفاوت خود یک مسئله مهندسی جدی است [12]. SWE-bench واحد ارزیابی را به درخواست مسئله در مخزن واقعی،
وابستگی‌ها و محیط اجرای آن گسترش می‌دهد [20]. این منابع برای جایگاه‌گذاری روش مهم‌اند، اما جمعیت کوچک
و جداگانه فاز ۱۳ را نماینده آن‌ها نمی‌کنند.

وظیفه سطح مخزن افزون بر مکان خطا، نسخه دقیق، درخت فایل، سامانه ساخت، وابستگی و فرمان ارزیابی را دربر
می‌گیرد. خرابی هرکدام ممکن است پیش از تولید وصله رخ دهد. بنابراین این مقاله مخرج برنامه‌ریزی‌شده،
فراخوانی‌های واقعی ارائه‌دهنده و نتایج علمی معتبر را جدا گزارش می‌کند؛ دو ردیف black-4 به دلیل شکست
checkout وارد نرخ موفقیت نمی‌شوند.

\#\#\# ۳.۳. عامل‌ها، بازیابی و گراف برنامه

ReAct چرخه استدلال و عمل [17] و Reflexion استفاده از بازخورد زبانی در تلاش‌های بعدی را صورت‌بندی
می‌کنند [18]. RepairAgent، AutoCodeRover و SpecRover این ایده‌ها را در تعمیر مخزن، جست‌وجوی ساختاری
و بازبینی مبتنی بر قصد پی می‌گیرند [21]، [22]، [23]. در ArchGuard-APR نام نقش‌ها به معنای چند مدل
مستقل نیست؛ نقش‌ها مرز داده و اختیار را تعریف می‌کنند و همه انتقال‌ها رکورد نوع‌دار دارند.

بازیابی افزوده، شواهد بیرونی را از دانش پارامتری جدا می‌سازد [15]. شبکه‌های عصبی گرافی برای کد [11]
و GraphCodeBERT [16] نشان می‌دهند رابطه ساختاری می‌تواند مفید باشد. بااین‌حال گراف این سامانه یک نمایش
یادگرفته‌شده نیست: از تصویر AST معتبر، با مرتب‌سازی قطعی و شناسه پایدار ساخته می‌شود و فقط برای
پرس‌وجوی قابل ممیزی به کار می‌رود.

\#\#\# ۳.۴. جایگاه رویکرد آگاه از معماری

معماری در این پژوهش یک پروفایل صریح و نسخه‌دار است، نه استنتاج احتمالی از نام فایل‌ها. مدل بازتابی
وابستگی میان مدل و پیاده‌سازی را قابل سنجش می‌کند [2] و تعریف کلاسیک معماری بر عناصر و روابط و منطق
طراحی تأکید دارد [1]. ArchGuard-APR این دیدگاه را به مرز تعمیر می‌آورد: گزارش معماری می‌تواند زمینه
تولید را غنی کند و پس از اعمال وصله نیز یک دروازه انطباق باشد. مقایسه نظام‌مند با کارهای پیشین در
[جدول ۱۰](../tables/table-related-work-comparison.md) آمده و صرفاً تفاوت دامنه و سازوکار را نشان می‌دهد،
نه برتری تجربی.

\#\# ۴. پرسش‌های پژوهش

**RQ1:** آیا قیود معماری، گراف مخزن، بازیابی زمینه و بازخورد اجرا را می‌توان در خط لوله‌ای بازتولیدپذیر
یکپارچه کرد؟ **RQ2:** چه محدودیت‌های عملیاتی در اجرای وظایف واجد شرایط پدیدار می‌شود؟ **RQ3:** چه
پیامدها، مصرف توکن و هزینه‌ای مشاهده می‌شود؟ هیچ پرسشی برتری را پیش‌فرض نمی‌گیرد.

[جدول ۲](../tables/table-2-research-questions.md) هر پرسش را به واحد تحلیل، سنجه و منبع شواهد پیوند
می‌دهد. RQ1 با مصنوعات و آزمون‌های یکپارچه‌سازی پاسخ توصیفی می‌گیرد؛ RQ2 مراحل توقف و ریزش را بررسی
می‌کند؛ و RQ3 فقط شمارش پیامد، توکن و هزینه مشاهده‌شده را گزارش می‌دهد. این تفکیک مانع آن است که
کامل‌بودن مهندسی به‌اشتباه شاهد اثربخشی علمی تلقی شود.

\#\# ۵. طراحی ArchGuard-APR

ورودی‌های خط فرمان، خدمات کاربرد را از راه قراردادهای دامنه فراخوانی می‌کنند. سازگارسازها عملیات Git،
ذخیره مصنوعات، ظرف اجرا و ارائه‌دهنده را پیاده می‌کنند. جریان داده یک‌طرفه است: مخزن به تصویر معتبر،
تصویر به گراف، گراف و پروفایل به گزارش معماری، و گزارش و پرس‌وجوی زمینه به بسته محدود تعمیر تبدیل
می‌شوند. هیچ مرحله پایین‌دستی برای تکمیل اطلاعات، کد هدف را دوباره نمی‌خواند.

![شکل ۱. معماری سامانه و مرزهای میان ورودی، هسته مستقل از فروشنده و سازگارسازهای زیرساختی.](../figures/system-architecture.mmd)

[جدول ۱](../tables/table-1-system-components.md) مؤلفه‌ها، ورودی و خروجی و مسئولیت هر مرز را خلاصه
می‌کند. قراردادهای دامنه، شناسه اجرای پایدار، تصویر مخزن، گراف، پروفایل کامپایل‌شده، نامزد وصله و
نتیجه ارزیابی را تعریف می‌کنند. لایه کاربرد ترتیب عملیات و سیاست توقف را نگه می‌دارد، اما Git، Docker
یا قالب یک ارائه‌دهنده را نمی‌شناسد. این جداسازی امکان آزمون کامل مسیر با مخزن و انتقال جعلی را فراهم
می‌کند و از ورود جزئیات فروشنده به تصمیم علمی جلوگیری می‌کند.

\#\#\# ۵.۱. مسئولیت، اعتماد و جریان داده

مرزهای اعتماد میان کد میزبان، مخزن هدف، ارائه‌دهنده و محیط اجرای ایزوله صریح‌اند. شناسه‌ها و اثرانگشت‌ها
در هر تبدیل ثبت می‌شوند. کنترل هزینه پیش از آغاز تماس ظرفیت لازم را رزرو می‌کند و پس از دریافت مصرف
واقعی، رزرو را آزاد و هزینه مشاهده‌شده را دقیقاً یک‌بار ثبت می‌کند. هیچ کلید یا پاسخ مدل در بسته مقاله
قرار ندارد.

اعتماد در چهار ناحیه توزیع شده است. مخزن هدف داده‌ای بالقوه خصمانه و فقط‌خواندنی است؛ میزبان مسئول
اعتبارسنجی و ایجاد مصنوعات است؛ محیط اجرا باید از میزبان و شبکه جدا باشد؛ و ارائه‌دهنده فقط بسته مجاز
را دریافت می‌کند. هر تبدیل پیش‌شرط خود را دوباره اعتبارسنجی می‌کند، اما منبع را برای بازسازی اطلاعات
پنهان دور نمی‌زند. شکست احراز هویت، بودجه، نشت، اثرانگشت یا پاک‌سازی به‌صورت بسته متوقف می‌شود.

\#\# ۶. تحلیل مخزن و ساخت گراف

تحلیل‌گر بدون واردکردن یا اجرای کد هدف، فایل‌های مجاز را با AST می‌خواند و تصویر استاندارد تولید
می‌کند. گراف شامل ماژول، تعریف، import و رابطه دربرگیری است. مرتب‌سازی پایدار و JSON استاندارد، مصنوعات
یکسان را در مکان‌های متفاوت همسان نگه می‌دارد. ریشه محافظت‌شده مانع نوشتن خروجی درون مخزن تحلیل‌شده
می‌شود.

تحلیل AST نام فایل، ماژول، کلاس، تابع و import را بدون importکردن بسته هدف استخراج می‌کند. رکوردها
پیش از هش‌گذاری بر اساس کلیدهای پایدار مرتب می‌شوند و مسیرها نسبت به ریشه معتبر ذخیره می‌شوند. ساخت
گراف فقط همین تصویر را می‌پذیرد؛ بنابراین تغییر فایل پس از تحلیل یا بازخوانی اتفاقی منبع نمی‌تواند
بی‌صدا ساختار ارزیابی‌شده را عوض کند. این انتخاب دامنه تحلیل را محدود می‌کند، ولی منشأ هر یال و گره
را قابل پیگیری نگه می‌دارد.

\#\#\# ۶.۱. ناوردای تصویر و گراف

گراف فقط از تصویر اعتبارسنجی‌شده ساخته می‌شود. مسیر خارج از ریشه، پیوند نمادین ناایمن، gitlink، فایل
تغییرکرده یا اثرانگشت ناسازگار باعث توقف می‌شود. این سخت‌گیری پوشش تحلیل پویا، reflection و کد تولیدی
را تضمین نمی‌کند؛ چنین مواردی در تهدیدهای اعتبار باقی می‌مانند.

[جدول ۳](../tables/table-3-qualified-population.md) جمعیت واجد شرایط منجمد را نشان می‌دهد. بازبینی
تغییرناپذیر، نشانی upstream، درخت Git، نبود آلودگی worktree و سیاست پیوندها پیش از هر تحلیل بررسی
می‌شود. هدف این کنترل‌ها حذف تفاوت‌های محیطی نیست، بلکه آشکارکردن آن‌ها پیش از آن است که به پیامد مدل
نسبت داده شوند.

\#\# ۷. مدل قید معماری

پروفایل YAML نسخه‌دار، قواعد انتساب و وابستگی را به بازنمایی استاندارد کامپایل می‌کند. چرخه، وابستگی
ممنوع و ماژول نامشخص با شناسه پایدار گزارش می‌شوند. پروفایل حاصل قضاوت انسانی است، نه حقیقت استخراج‌شده
خودکار؛ ازاین‌رو منشأ، نسخه و بازبینی آن بخشی از شواهد است.

\#\#\# ۷.۱. چرخه عمر پروفایل و تفسیر قواعد

بارگذار YAML فقط نحو و طرح داده را می‌سنجد؛ کامپایلر انتخاب‌گرها را به بازنمایی استاندارد و قواعد
نسخه‌دار تبدیل می‌کند؛ و ارزیاب، پروفایل کامپایل‌شده را روی گراف معتبر اجرا می‌کند. جدایی این سه مرحله
مانع آن است که خطای بارگذاری با نقض معماری یا خطای ذخیره‌سازی مخلوط شود. گزارش شامل قواعد فعال، قواعد
ردشده، ماژول‌های منتسب و نامنتسب و اثرانگشت ورودی‌هاست.

قواعد وابستگی می‌توانند مبدأ، مقصد و شدت را مشخص کنند. ارزیابی چرخه و وابستگی ممنوع قطعی است و مسیر
نقض را با شناسه‌های منبع‌آزاد گزارش می‌کند. ماژول مبهم، انتخاب‌گر هم‌پوشان یا انتساب ناقص به‌صورت
محافظه‌کارانه پذیرفته نمی‌شود. این مدل جایگزین قضاوت معمار نیست؛ تصمیم انسانی را به قراردادی قابل
نسخه‌بندی و آزمون تبدیل می‌کند.

\#\# ۸. گردش کار تعمیر چندعاملی

هماهنگ‌کننده، تحلیل، بازیابی، تولید، اعتبارسنجی و ارزیابی را با حالت‌های نوع‌دار پیش می‌برد. نقش‌ها
مسئولیت را جدا می‌کنند، اما به معنای استقلال علمی چند مدل نیستند. تماس ارائه‌دهنده تنها پس از تأیید
پیکربندی منجمد، پذیرش هزینه و گذر از کنترل امنیتی ممکن است. وصله نامعتبر به نتیجه پایانی صریح تبدیل
می‌شود و با تصحیح خاموش قالب، موفق جلوه داده نمی‌شود.

![شکل ۲. گردش کار تعمیر از آماده‌سازی قطعی تا تولید، اعتبارسنجی، اجرای ایزوله و ثبت شواهد.](../figures/repair-workflow.mmd)

\#\#\# ۸.۱. چرخه نامزد و حسابداری

هر تماس شناسه یکتا، وضعیت شروع و پایان، مصرف و هزینه دارد. تلاش مجدد در آزمایش نهایی صفر بود. شمارش
فراخوانی و هزینه به شرط تکمیل مصرف دقیقاً یک‌بار انجام می‌شود. طبقه‌بندی زیرساختی پیش از ارائه‌دهنده
از طبقه‌بندی علمی پس از پاسخ جداست.

یک درخواست تولید، شناسه وظیفه، شرط، اثرانگشت manifest زمینه، پیکربندی ارائه‌دهنده و رزرو بودجه را
به هم می‌بندد. پس از پاسخ، تجزیه‌کننده فقط پوشش وصله منجمد را می‌پذیرد؛ مسیر ناامن، فایل ممنوع، اندازه
بیش از سقف یا نحو مخدوش به‌طور دستی اصلاح نمی‌شود. نامزد مجاز در worktree موقت اعمال و با فرمان منجمد
ارزیابی می‌شود. پاک‌سازی تعهد پایانی است و نقص آن حتی پس از نتیجه ظاهراً موفق، توقف fatal محسوب می‌شود.

\#\#\# ۸.۲. نقش‌ها و انتزاع ارائه‌دهنده

تحلیل، بازیابی، تولید، اعتبارسنجی و ارزیابی ورودی و سطح اعتماد متفاوت دارند و از طریق رکوردهای نوع‌دار
ارتباط برقرار می‌کنند، نه یک گفت‌وگوی مشترک تغییرپذیر. انتقال جعلی در آزمون‌های آفلاین فراخوانی را ثبت
و پاسخ fixture را بازمی‌گرداند. تنها سازگارساز زنده اعتبارنامه و قالب سیمی فروشنده را تفسیر می‌کند.
در نتیجه، منطق توقف، بودجه و طبقه‌بندی در کاربرد باقی می‌ماند و با تعویض ارائه‌دهنده تغییر نمی‌کند.

\#\# ۹. بازیابی زمینه و بازخورد اجرا

پرس‌وجوی زمینه همسایگی ماژول و نماد را با سقف تعداد و اندازه انتخاب می‌کند و منشأ هر قطعه را در manifest
می‌نویسد. اجرای وصله پس از کنترل مسیر و محتوا در محیط محدود انجام می‌شود. بازخورد شامل وضعیت‌های
استاندارد و اثرانگشت است و متن حساس، کد منبع، prompt یا پاسخ خام را منتشر نمی‌کند.

\#\#\# ۹.۱. بودجه بازیابی و منشأ

پرس‌وجو با ماژول یا نماد آغاز می‌شود و همسایه‌ها را در ترتیب ثابت تا سقف تعداد گره، یال و بایت طی
می‌کند. انتخاب هر قطعه با شناسه تصویر، فاصله گراف، دلیل انتخاب و اندازه ثبت می‌شود. اگر یک قطعه سقف را
پر کند، سامانه آن را بی‌صدا نصف یا با فایل دیگری جایگزین نمی‌کند؛ تصمیم حذف در manifest آشکار می‌ماند.
این رفتار مقایسه دو شرط را ممکن می‌کند و از رشد نامحدود prompt جلوگیری می‌کند.

زمینه معماری شامل قواعد مرتبط و نقض‌های موجود است، نه کل پروفایل. شرط خط پایه همان مسئله، نسخه، مدل،
ابزار و سیاست وصله را دارد، اما این بسته معماری را دریافت نمی‌کند. تفاوت تیمار می‌تواند طول ورودی را
تغییر دهد؛ به همین دلیل توکن و هزینه نتیجه ثانویه صریح‌اند و نباید با موفقیت تعمیر یکی گرفته شوند.

\#\#\# ۹.۲. بازخورد ایزوله

پس از اعمال نامزد، ارزیاب در محیط محدود با سقف زمان، پردازنده، حافظه، فرایند و خروجی اجرا می‌شود.
شبکه و mountها مطابق پروتکل بسته‌اند و تنها وضعیت استاندارد، شمارش و اثرانگشت از مرز خارج می‌شود.
خروجی خام ممکن است حاوی کد، مسیر محلی یا داده حساس باشد و در بسته انتشار قرار نمی‌گیرد. بازخورد برای
حلقه بعدی نیز فقط از درگاه پالایش‌شده عبور می‌کند.

این کنترل‌ها اجرای کد نامطمئن را بی‌خطر مطلق نمی‌کنند. آن‌ها سطح دسترسی و دامنه خسارت را کاهش می‌دهند
و شکست مهار را قابل تشخیص می‌سازند. هر خطای آماده‌سازی پیش از تماس، از جمله checkout، نتیجه علمی نیست؛
هر پاسخ دریافت‌شده نیز تا گذر از قرارداد وصله، تعمیر محسوب نمی‌شود.

\#\# ۱۰. امنیت، ایزوله‌سازی و بازتولیدپذیری

خط لوله در برابر خروج از مسیر، نشت پاسخ طلایی، دسترسی شبکه نامجاز و نوشتن در مخزن هدف بسته عمل می‌کند.
سقف CPU، حافظه، فرایند، زمان و خروجی، دامنه خسارت را کاهش می‌دهد ولی خطر را صفر نمی‌کند. شواهد با
ترتیب پایدار، زمان‌های حذف‌شده از محتوای استاندارد و اثرانگشت SHA-256 ثبت می‌شوند.

زنجیره امنیت از کسب مخزن آغاز می‌شود: بازبینی و remote باید با جمعیت منجمد منطبق باشند، درخت نباید
gitlink یا پیوند خارج‌شونده غیرمجاز داشته باشد و worktree باید پاک بماند. تحلیل هرگز کد هدف را import
یا اجرا نمی‌کند. خروجی نیز خارج از ریشه محافظت‌شده نوشته می‌شود. در مرحله وصله، canonicalization مسیر
پیش از اعمال سیاست انجام می‌شود تا شکل‌های متفاوت یک مسیر نتوانند کنترل را دور بزنند.

بازتولیدپذیری در اینجا به معنای بازسازی قطعی مصنوعات استاندارد از ورودی منجمد است، نه تضمین پاسخ
بایت‌به‌بایت یک مدل احتمالی. نسخه ابزار، seedهای موجود، ترتیب اجرا، اثرانگشت پیکربندی و شناسه تماس ثبت
می‌شوند. مصنوعات علمی، منبع هدف، prompt، پاسخ خام و راز را حمل نمی‌کنند. این مرزبندی امکان ممیزی عمومی
آمار را فراهم می‌کند، درحالی‌که محدودیت محرمانگی و ایمنی حفظ می‌شود.

\#\# ۱۱. روش‌شناسی آزمایش

پروتکل اکتشافی سه وظیفه را در دو شرط خط پایه و ArchGuard، در مجموع شش اجرا، برنامه‌ریزی کرد. ترتیب،
جمعیت، مدل، پیکربندی ارائه‌دهنده و بودجه پیش از مشاهده نتیجه منجمد شدند. معیار موفقیت پروژه‌محور و
سیاست وصله سخت‌گیرانه بود.

![شکل ۳. پروتکل زوجی منجمد با سه وظیفه و دو شرط برای هر وظیفه.](../figures/experimental-protocol.mmd)

\#\#\# ۱۱.۱. احراز صلاحیت بدون مشاهده پیامد

وظایف بر اساس مجوز، دسترسی به بازبینی تغییرناپذیر، قابلیت اجرای ارزیاب و محدودیت منابع انتخاب شدند.
این فرایند از وصله طلایی برای هدایت تولید استفاده نکرد. تنها دو زوج PySnooper نتیجه علمی کامل ساختند؛
یک پروژه در نتایج معتبر نمایندگی شد.

ترتیب جمعیت بخشی از پروتکل بود: هر وظیفه ابتدا جایگاه خط پایه و سپس جایگاه ArchGuard داشت. ordinalها
در checkpoint پایدار ماندند تا ادامه اجرا شواهد کامل‌شده را تکرار نکند. صلاحیت بر مجوز، دسترسی نسخه،
فرمان ارزیاب و منابع تکیه داشت و نتیجه تعمیر یا وصله طلایی را مشاهده نکرد. این فرایند سوگیری ناشی از
انتخاب پس از مشاهده نتیجه را کاهش می‌دهد، اما سه وظیفه را نماینده کل APR نمی‌سازد.

\#\#\# ۱۱.۲. ساخت تیمار و انجماد

خط پایه زمینه مسئله و مخزن را بدون گزارش معماری دریافت کرد. شرط ArchGuard بسته زمینه و بازخورد قیود
معماری را نیز دریافت کرد. هر دو از مدل `gpt-5.4-2026-03-05` و سقف‌های یکسان ارائه‌دهنده استفاده کردند.
این گزارش هیچ اجرای تازه‌ای ایجاد نمی‌کند و تنها شواهد منجمد را بازگو می‌کند.

قیمت‌ها و سقف ۴٫۵۰ دلار پیش از اجرا تثبیت شد. آغاز هر تماس به تأیید صریح هزینه و پردازش منبع نیاز
داشت. runner بدترین هزینه مجاز را پیش از ارسال رزرو و پس از دریافت usage با هزینه واقعی مصالحه کرد.
جمعیت، پروتکل و پیکربندی ارائه‌دهنده در آغاز و resume با اثرانگشت کنترل شدند تا ادامه اجرا نتواند شرط
یا مدل را عوض کند.

\#\#\# ۱۱.۳. سیاست طبقه‌بندی و ریزش

چهار پاسخ دریافت‌شده به دلیل عبور نکردن از قرارداد وصله، `invalid_patch` شدند. دو اجرای black-4 در
مرحله checkout و با رده `validation_failure` پیش از تماس ارائه‌دهنده متوقف شدند. اثرانگشت تشخیص
`71f60e3969313f139c493cad31ee20640fd6c2870d70411debdd5ab1956994e1` است. این دو مورد شکست مدل یا
تعمیر نیستند.

[جدول ۷](../tables/table-7-attrition.md) مخرج شش جایگاه و دو ریزش را جداگانه ثبت می‌کند. ordinal پنج
در وظیفه black-4 و شرط خط پایه هنگام اعتبارسنجی checkout متوقف شد؛ چون آماده‌سازی منصفانه وظیفه ممکن
نبود، ordinal شش نیز به‌عنوان ریزش زوجی بسته شد. هیچ‌کدام طبقه وصله، مصرف توکن، هزینه یا پیامد
ارائه‌دهنده ندارند.

\#\#\# ۱۱.۴. کنترل‌های بازتولیدپذیری

اثرانگشت جمعیت `421f3129ef6bee5883c41d05f087d6faff5b2a3de3b2ea543b4d27ade1488886`، پروتکل
`eae1920ce71c025d2d94096798daf8b14df098249743631efe9a0aa7a541fbe8` و پیکربندی ارائه‌دهنده
`f5dc6fb5b8bf54c7cf1b58b990fa3db44073500df9903e3a7d93cbeb1ad92c2d` است. ثبت نسخه‌ها، ترتیب
اجرا و هزینه، ممیزی را ممکن می‌کند؛ بازتولید دقیق رفتار احتمالی ارائه‌دهنده تضمین نمی‌شود.

checkpoint اجراهای حمل‌شده، شناسه تماس، usage، هزینه و حالت پاک‌سازی را به هم می‌بندد. سیاست resume
میان علت توقف تاریخی و توقف فعال ادامه جدید فرق می‌گذارد و ordinal کامل‌شده را دوباره زمان‌بندی
نمی‌کند. واردسازی نهایی، SHA-256 مصنوع
`48080661ddbff3e486d47cebc63754e29d4dd4664febbb7f8143484e4c03e99d` را اعتبارسنجی کرد. نمایه
کنترل‌های قابل بازسازی در [جدول ۹](../tables/table-9-reproducibility.md) آمده است.

\#\# ۱۲. نتایج اکتشافی

[جدول ۴](../tables/table-4-per-execution-results.md) چهار اجرای کامل و دو ردیف ریزش را بدون ادغام
مخرج‌ها نشان می‌دهد و [جدول ۵](../tables/table-5-paired-comparison.md) مقایسه را به دو وظیفه زوجی معتبر
محدود می‌کند. **خلاصه نتایج و هزینه:** خط پایه: ۲ زوج معتبر، موفقیت ۰/۲، ۲ وصله نامعتبر،
۷٬۸۲۱ توکن ورودی، ۱٬۰۳۷ توکن خروجی و هزینه ۰٫۰۳۵۱۰۷۵ دلار. ArchGuard: ۲ زوج معتبر،
موفقیت ۰/۲، ۲ وصله نامعتبر، ۲۷٬۶۱۰ توکن ورودی، ۱٬۷۰۰ توکن خروجی و هزینه ۰٫۰۹۴۵۲۵ دلار.
مجموع: ۲ زوج معتبر، موفقیت ۰/۴، ۴ وصله نامعتبر، ۳۵٬۴۳۱ توکن ورودی، ۲٬۷۳۷ توکن
خروجی و هزینه ۰٫۱۲۹۶۳۲۵ دلار.

![شکل ۴. حسابداری پیامدها: چهار وصله نامعتبر و دو ریزش زیرساختی.](../figures/execution-outcomes.svg)

از شش اجرای برنامه‌ریزی‌شده، چهار اجرا کامل و دو اجرا دچار ریزش شدند؛ نرخ ریزش زیرساختی ۳۳٫۳ درصد
است. هر دو شرط ۰ از ۲ موفقیت داشتند و هر چهار خروجی کامل‌شده وصله نامعتبر بود. بنابراین تفاوت اثربخشی
مشاهده نشد، اما برابری دو روش نیز اثبات نمی‌شود. شرط ArchGuard در این نمونه کوچک ورودی و هزینه بیشتری
مصرف کرد؛ این مشاهده قابل تعمیم نیست.

[جدول ۶](../tables/table-6-token-cost.md) usage و هزینه را به تفکیک شرط نگه می‌دارد. چهار فراخوانی
تکمیل شد، retry برابر صفر و هزینه رزروشده نهایی صفر بود؛ ۴٫۳۷۰۳۶۷۵ دلار از سقف باقی ماند.

![شکل ۵. هزینه مشاهده‌شده ارائه‌دهنده به تفکیک شرط در نمونه اکتشافی.](../figures/cost-comparison.svg)

![شکل ۶. توکن ورودی مشاهده‌شده به تفکیک شرط.](../figures/input-token-comparison.svg)

\#\#\# ۱۲.۱. تفسیر توصیفی

تنها دو زوج معتبر و صفر زوج ناهمخوان موفق وجود دارد. آزمون معناداری یا برآورد اثر پایدار با این داده
موجه نیست. این ارزیابی باید امکان‌پذیری اکتشافی تلقی شود، نه ارزیابی تأییدی اثربخشی.

مخرج درست به پرسش بستگی دارد. تکمیل عملیاتی ۴ از ۶ جایگاه و ریزش ۲ از ۶ یا ۳۳٫۳ درصد است. اثربخشی
تعمیر روی دو وظیفه زوجی معتبر گزارش می‌شود: خط پایه ۰/۲ و ArchGuard نیز ۰/۲. بازده قرارداد وصله صفر
از چهار پاسخ کامل است و حسابداری ارائه‌دهنده چهار تماس دارای usage را نشان می‌دهد. جایگزین‌کردن یکی
از این مخرج‌ها با دیگری، هم ریزش را پنهان می‌کند و هم نتیجه مدل را مخدوش می‌سازد.

افزایش ورودی ArchGuard، ۲۷٬۶۱۰ در برابر ۷٬۸۲۱ توکن، با افزودن زمینه معماری سازگار است؛ خروجی‌ها
۱٬۷۰۰ در برابر ۱٬۰۳۷ توکن بودند. با چهار تماس نمی‌توان سهم سربار ثابت تیمار را از رفتار خاص وظیفه
جدا کرد. بنابراین این اعداد فقط توصیف تماس‌های منجمد با قیمت منجمدند و پیش‌بینی هزینه آینده نیستند.

\#\# ۱۳. تحلیل شکست

شواهد ثابت می‌کنند چهار خروجی در مرز سیاست وصله رد شدند، اما علت یگانه‌ای را اثبات نمی‌کنند. قالب،
دستور زبان وصله، سیاست مسیر، کفایت زمینه، یکپارچه‌سازی و رفتار مدل همگی سازوکارهای محتمل‌اند. نسبت دادن
قطعی شکست به یکی از آن‌ها بدون شواهد بیشتر نارواست.

\#\#\# ۱۳.۱. مرز وصله نامعتبر

اعتبارسنجی سخت‌گیرانه، ایمنی و وضوح شواهد را افزایش می‌دهد و هم‌زمان می‌تواند بازده ظاهری را کاهش دهد.
پژوهش بعدی باید تشخیص منبع‌آزاد دقیق‌تری برای مراحل تجزیه و سیاست فراهم کند، بدون آنکه پاسخ خام یا
جزئیات حساس را منتشر سازد.

\#\#\# ۱۳.۲. ریزش checkout

هر دو وضعیت black-4 پیش از تماس با ارائه‌دهنده در اعتبارسنجی checkout متوقف شدند. آن‌ها در مخرج نتایج
تعمیر قرار نمی‌گیرند و تنها به‌عنوان ریزش زیرساختی گزارش می‌شوند. این رخداد نیاز به آزمون بیشتر کسب
مخزن و ناوردای بازبینی را نشان می‌دهد، نه ضعف مدل را.

![شکل ۷. حسابداری ریزش، مخرج برنامه‌ریزی‌شده را حفظ و شکست پیش از ارائه‌دهنده را از پیامد تعمیر جدا می‌کند.](../figures/attrition-accounting.mmd)

\#\# ۱۴. بحث

پیاده‌سازی نشان می‌دهد می‌توان قیود معماری را با خط لوله‌ای قابل ممیزی ترکیب کرد. بااین‌حال، داده‌ها
هیچ برتری تعمیر را نشان نمی‌دهند. RQ1 از نظر یکپارچه‌سازی مهندسی پشتیبانی می‌شود؛ RQ2 اعتبار وصله و
پایداری زیرساخت را محدودیت‌های اصلی می‌یابد؛ و RQ3 چهار وصله نامعتبر، ۳۸٬۱۶۸ توکن و هزینه
۰٫۱۲۹۶۳۲۵ دلار را ثبت می‌کند.

افزایش هزینه شرط ArchGuard با زمینه ورودی بزرگ‌تر سازگار است، اما دو زوج برای نتیجه‌گیری درباره نسبت
هزینه و فایده کافی نیست. ارزش گراف و قید نیز از نرخ موفقیت صفر قابل جداکردن نیست. مطالعه بعدی باید پس
از رفع مشکلات اعتبار وصله و checkout، جمعیت چندپروژه‌ای بزرگ‌تری را به‌صورت آینده‌نگر طراحی کند.

\#\#\# ۱۴.۱. آنچه پیاده‌سازی اثبات می‌کند

مخزن و آزمون‌های قطعی نشان می‌دهند پروفایل معماری، گراف معتبر، بازیابی بودجه‌دار، هماهنگی مستقل از
فروشنده، بازرسی وصله، ارزیابی ایزوله، حسابداری هزینه، checkpoint ادامه‌پذیر و شواهد استاندارد را می‌توان
پشت درگاه‌های صریح به هم متصل کرد. همچنین مسیر ارائه‌دهنده را می‌توان با انتقال جعلی و بدون شبکه آزمود
و آمار مقاله را از داده‌های منجمد بازساخت.

این شواهد نشان نمی‌دهند زمینه معماری برای مدل مفید بوده، این گراف بهترین نمایش است، یا چند نقش از یک
حلقه واحد بهترند. پوشش آزمون و کامل‌بودن مهندسی درباره انضباط پیاده‌سازی‌اند، نه کیفیت وصله. چهار خروجی
نامعتبر نیز برآوردی پایدار از اثربخشی تعمیر فراهم نمی‌کنند.

\#\#\# ۱۴.۲. زمینه معماری و هزینه

اطلاعات معماری از بودجه زمینه مصرف می‌کند. ممکن است مولد را از نقض وابستگی دور کند، یا شواهد محلی
مهم‌تر را کنار بزند و قواعد نامرتبط اضافه کند. سیاست بازیابی باید اقلام انتخاب‌شده، منشأ، دلیل و سهم
بودجه آن‌ها را آشکار کند. مطالعه آینده باید ترکیب زمینه را پیشاپیش ثبت و اثر آن را بر قابلیت پذیرش،
موفقیت رفتاری و انطباق معماری جداگانه بسنجد.

در داده حاضر ArchGuard ورودی و هزینه بیشتری داشت، اما هیچ شرطی وصله معتبر نساخت. در نتیجه نسبت هزینه
به فایده قابل محاسبه نیست. طرح بعدی باید هزینه به ازای تلاش، وصله معتبر، وصله plausible و وصله صحیح
را از پیش تعریف کند و قیمت و مدل را همراه نتیجه گزارش دهد.

\#\#\# ۱۴.۳. درس‌هایی برای مطالعات بازتولیدپذیر APR

نخست، طبقه‌های پیامد باید همه خط لوله را پوشش دهند، نه فقط عبور یا شکست آزمون. دوم، مطالعه پولی
ادامه‌پذیر به شمارش دقیقاً یک‌باره تماس و هزینه نیاز دارد. سوم، احراز صلاحیت باید کسب، محیط و پاک‌سازی
را بدون نگاه به نتیجه تعمیر بیازماید. چهارم، خروجی منفی و ریزش نباید حذف شوند. در نهایت، ادعا باید از
مخرج تحقق‌یافته پیروی کند: برنامه شش‌تایی با دو توقف پیش از ارائه‌دهنده، شش اجرای کامل نیست.

\#\# ۱۵. تهدیدهای اعتبار

اعتبار سازه به تعریف سخت‌گیرانه و پروژه‌محور موفقیت وابسته است. اعتبار درونی با علت نامعلوم وصله‌های
نامعتبر و وابستگی محیطی محدود می‌شود. اعتبار بیرونی تنها بر دو زوج از یک پروژه تکیه دارد. اعتبار
نتیجه‌گیری با نبود تعمیر موفق و صفر زوج ناهمخوان، هرگونه استنباط یا ادعای برتری را رد می‌کند. ریزش
زیرساختی ۳۳٫۳ درصد است. نتایج به مدل، ارائه‌دهنده، زمان و پیکربندی وابسته‌اند؛ تحلیل ایستا رفتار پویا
را کامل نمی‌بیند؛ پروفایل به قضاوت انسانی نیاز دارد؛ و ایزوله‌سازی همه خطرها را حذف نمی‌کند.

**اعتبار سازه.** موفقیت در این مطالعه فقط عبور از آزمون نیست: وصله باید قابل تجزیه، مطابق سیاست،
قابل اعمال، منطبق با معماری و در ارزیاب موفق باشد. این تعریف آگاهانه سخت‌گیرانه است و با پژوهشی که
متن پاسخ را دستی اصلاح یا تنها test adequacy را می‌سنجد یکسان نیست. طبقه `invalid_patch` مرزی قطعی
اما درشت‌دانه است؛ قابلیت بازتولید را حفظ می‌کند و در عوض تفکیک علت نحوی، envelope، مسیر یا اندازه را
محدود می‌سازد.

**اعتبار درونی.** انجماد ترتیب و اثرانگشت‌ها drift پنهان را کاهش می‌دهد، ولی رفتار احتمالی مدل، ساخت
prompt، کد یکپارچه‌سازی و وضعیت محیط می‌توانند پیامد را تغییر دهند. چون پاسخ خام عمداً در مصنوع
منبع‌آزاد نیست، داده موجود علت چهار وصله نامعتبر را از میان این سازوکارها جدا نمی‌کند. نسبت‌دادن همه
شکست‌ها به مدل، parser یا زمینه معماری فراتر از شواهد است.

**اعتبار بیرونی و انتخاب وظیفه.** فقط دو زوج معتبر از یک پروژه نتیجه علمی دارند. جمعیت اولیه سه وظیفه
بود، اما از دست‌رفتن black-4 تنوع پروژه را بیشتر کاهش داد. هیچ تعمیمی به زبان، نوع نقص، مخزن، مدل،
ارائه‌دهنده یا بنچمارک‌های Defects4J [6]، BugsInPy [14] و SWE-bench [20] پشتیبانی نمی‌شود. احراز
صلاحیت outcome-blind از انتخاب نتیجه‌محور جلوگیری می‌کند، نه اینکه نمونه را نماینده سازد.

**اعتبار نتیجه‌گیری.** هیچ تعمیر موفق و هیچ زوج موفق ناهمخوان وجود ندارد. اندازه نمونه برای آزمون
استنباطی، فاصله اثر پایدار، برابری یا برتری توان کافی ندارد. صفر بودن اختلاف مشاهده‌شده صرفاً حساب
۰/۲ در برابر ۰/۲ است. هزینه بیشتر ArchGuard نیز مشاهده چهار تماس است و اثر علّی یا پیش‌بینی اقتصادی
محسوب نمی‌شود.

**اعتبار زیرساخت.** دو جایگاه از شش جایگاه، یعنی ۳۳٫۳ درصد، هنگام checkout ریزش کردند. جداکردن آن‌ها
از پیامد تعمیر از آلودگی نتیجه جلوگیری می‌کند، اما مقایسه موردنظر black-4 را نامشاهده می‌گذارد. aging
وابستگی، دسترسی بازبینی تغییرناپذیر و تفاوت میزبان می‌توانند قابلیت اجرای وظایف تاریخی را تغییر دهند.
مطالعه بعدی باید preflight منبع‌آزاد همه وظایف را شرط freeze کند.

**نمایش مخزن و پروفایل انسانی.** AST ایستا ممکن است import پویا، reflection، کد تولیدی یا رابطه زمان
اجرا را نبیند. پروفایل معماری نیز قضاوت بازبین را کدگذاری می‌کند و می‌تواند ناقص یا محل اختلاف باشد.
ارزیابی آینده باید پوشش نمایش، دلیل قواعد و توافق میان بازبینان را بسنجد؛ سامانه حاضر معماری را حقیقتی
خودکار و بی‌خطا معرفی نمی‌کند.

**وابستگی مدل، ارائه‌دهنده و امنیت.** نتایج به `gpt-5.4-2026-03-05`، قیمت و پیکربندی منجمد و زمان
اجرا وابسته‌اند. sandbox، کنترل شبکه و سیاست مسیر خطر را کاهش می‌دهند، ولی آسیب‌پذیری runtime، وابستگی
یا میزبان را حذف نمی‌کنند. بازتولید مصنوعات و حساب‌ها به معنای بازتولید بایت‌به‌بایت پاسخ احتمالی
ارائه‌دهنده نیست. جمع‌بندی ساختاریافته این تهدیدها در [جدول ۸](../tables/table-8-threats.md) است.

\#\# ۱۶. ملاحظات اخلاقی و بازتولیدپذیری

نتایج منفی و هزینه‌ها حفظ شده‌اند. تنها منابع تأییدشده به کار رفته‌اند. هیچ کد هدف، وصله طلایی، کلید،
prompt یا پاسخ ارائه‌دهنده در بسته مقاله نیست. مسئولیت تفسیر، افشا و ارسال بر عهده نویسنده انسانی
است. اثرانگشت نتیجه نهایی `761d9853f6ecfb4dafdaad81de781c4d49108b8b82912aeec509e173c7d56c70`
است.

گزارش نتایج منفی از گزینش مطلوب جلوگیری می‌کند و ثبت هزینه، اثر محیطی و اقتصادی تماس‌ها را قابل
بررسی می‌سازد. در مقابل، انتشار کد هدف، پاسخ خام یا وصله طلایی می‌توانست محرمانگی، مجوز یا اعتبار
ارزیابی آینده را تهدید کند؛ بنابراین بسته تنها رکوردهای استاندارد لازم برای بازشماری ادعاها را نگه
می‌دارد. نویسنده انسانی مسئول انتخاب ادعا، ترجمه، افشا، بررسی ارجاع و تصمیم نهایی انتشار است.

[جدول ۹](../tables/table-9-reproducibility.md) مسیر داده‌های نتایج، جدول ادعا-شواهد، اثرانگشت‌ها و
آزمون‌های قطعی را فهرست می‌کند. هر عدد در متن باید از ردیف‌های canonical بازساخته شود و prose منبع
حقیقت نیست. هر اصلاح آینده باید با منشأ و نسخه جدید ثبت شود و نتیجه منجمد فاز ۱۳ را بازنویسی نکند.

\#\# ۱۷. نتیجه‌گیری و کار آینده

ArchGuard-APR قیود معماری، گراف مخزن، بازیابی زمینه، هماهنگی نقش‌محور، اعتبارسنجی وصله و بازخورد
ایزوله را در زیرساختی قابل ممیزی گرد هم می‌آورد. فاز ۱۳ امکان ترکیب این اجزا را نشان داد، ولی شاهدی
برای مزیت تعمیر فراهم نکرد. اولویت آینده، تولید وصله معتبر، تشخیص بهتر، پایداری checkout و سپس ارزیابی
آینده‌نگر بزرگ‌تر است.

کار آینده باید به ترتیب وابستگی پیش رود. ابتدا validator باید رده‌های منبع‌آزاد از پیش تعریف‌شده برای
envelope، نحو، مسیر، اندازه و اعمال تولید کند و انتقال‌های جعلی، وصله‌های پذیرفتنی را تا ارزیاب دنبال
کنند. سپس کسب مخزن باید روی همه وظایف منجمد بدون تضعیف remote، revision، tree، symlink، gitlink،
پاکی worktree و نشت تمرین شود. تنها پس از رفع این دو مانع، پروتکل چندپروژه‌ای با توان و معیارهای هزینه
از پیش تعیین‌شده موجه است.

یک ارزیابی بزرگ‌تر باید ترکیب زمینه را نیز اندازه‌گیری کند: چه مقدار از prompt مسئله، کد محلی، رابطه
گراف و قاعده معماری است و هر جزء چرا انتخاب شده است. پیامدها باید پذیرش وصله، آزمون رفتاری، انطباق
معماری، هزینه و زمان را جدا نگه دارند. تا آن زمان، نتیجه این مقاله امکان‌پذیری مهندسی و محدودیت‌های
مشاهده‌شده است، نه ادعای اثربخشی، برتری یا تعمیم.

\#\# دسترس‌پذیری داده‌ها

داده‌های استاندارد و عاری از کد منبع در `research/results/phase-13/` و داده‌های انتشار در `paper/data/`
قرار دارند. هیچ کد هدف یا وصله طلایی در این بسته توزیع نشده است.

\#\# دسترس‌پذیری مصنوعات

مخزن شامل جدول‌ها، شکل‌ها، ماتریس ادعا-شواهد، اثرانگشت‌ها، اسکریپت ساخت قطعی و آزمون‌های راستی‌آزمایی
است. مقاله ارسال یا پذیرفته نشده و محل انتشار هنوز تعیین نشده است. محدودیت عدم انتشار prompt و پاسخ
خام آگاهانه و برای حفظ ایمنی و منشأ داده است.

\#\# مشارکت نویسنده

نویسنده به‌تنهایی مسئول مفهوم‌پردازی، طراحی روش پژوهش، توسعه نرم‌افزار، اعتبارسنجی، انجام بررسی‌ها،
مدیریت داده‌ها، تحلیل رسمی، مصورسازی و آماده‌سازی مقاله بوده است.

\#\# حمایت مالی

این پژوهش هیچ‌گونه حمایت مالی خارجی دریافت نکرده است.

\#\# تعارض منافع

نویسنده اعلام می‌کند که هیچ‌گونه تعارض منافعی وجود ندارد.

\#\# منابع

[1] D. E. Perry and A. L. Wolf, “Foundations for the Study of Software Architecture,” ACM SIGSOFT Software Engineering Notes, 1992, doi:10.1145/141874.141884.

[2] G. C. Murphy, D. Notkin, and K. J. Sullivan, “Software Reflexion Models: Bridging the Gap Between Source and High-Level Models,” ACM SIGSOFT FSE, 1995, doi:10.1145/222132.222136.

[3] M. Monperrus, “Automatic Software Repair: A Bibliography,” ACM Computing Surveys, 2018, doi:10.1145/3105906.

[4] C. Le Goues, T. Nguyen, S. Forrest, and W. Weimer, “GenProg: A Generic Method for Automatic Software Repair,” IEEE TSE, 2012, doi:10.1109/TSE.2011.104.

[5] H. D. T. Nguyen, D. Qi, A. Roychoudhury, and S. Chandra, “SemFix: Program Repair via Semantic Analysis,” ICSE, 2013, doi:10.1109/ICSE.2013.6606623.

[6] R. Just, D. Jalali, and M. D. Ernst, “Defects4J,” ISSTA, 2014, doi:10.1145/2610384.2628055.

[7] S. Mechtaev, J. Yi, and A. Roychoudhury, “Angelix,” ICSE, 2016, doi:10.1145/2884781.2884807.

[8] F. Long and M. Rinard, “Automatic Patch Generation by Learning Correct Code,” POPL, 2016, doi:10.1145/2837614.2837617.

[9] M. Martinez et al., “Automatic Repair of Real Bugs in Java,” Empirical Software Engineering, 2017, doi:10.1007/s10664-016-9470-4.

[10] R. K. Saha et al., “Elixir,” ASE, 2017, doi:10.1109/ASE.2017.8115675.

[11] M. Allamanis et al., “Learning to Represent Programs with Graphs,” ICLR, 2018.

[12] T. Durieux et al., “Empirical Review of Java Program Repair Tools,” ESEC/FSE, 2019, doi:10.1145/3338906.3338911.

[13] T. Lutellier et al., “CoCoNuT,” ISSTA, 2020, doi:10.1145/3395363.3397369.

[14] R. Widyasari et al., “BugsInPy,” ESEC/FSE, 2020, doi:10.1145/3368089.3417943.

[15] P. Lewis et al., “Retrieval-Augmented Generation for Knowledge-Intensive NLP Tasks,” NeurIPS, 2020.

[16] D. Guo et al., “GraphCodeBERT,” ICLR, 2021.

[17] S. Yao et al., “ReAct,” ICLR, 2023.

[18] N. Shinn et al., “Reflexion,” NeurIPS, 2023, doi:10.52202/075280-0377.

[19] C. S. Xia et al., “Automated Program Repair in the Era of Large Pre-trained Language Models,” ICSE, 2023, doi:10.1109/ICSE48619.2023.00129.

[20] C. E. Jimenez et al., “SWE-bench,” ICLR, 2024.

[21] I. Bouzenia et al., “RepairAgent,” arXiv:2403.17134, 2024.

[22] Y. Zhang et al., “AutoCodeRover,” arXiv:2404.05427, 2024.

[23] H. Ruan et al., “SpecRover,” arXiv:2408.02232, 2024.

\#\# پیوست الف. اثرانگشت‌های بازتولیدپذیری

اثر انگشت manifest نهایی `8afb87add96e29898ae712f8ef08928b9ba1290a0aabc3ce34bac82adb4b86c2`
و اثر انگشت import برابر `80eac046fc614cad1beb7de127ba69aa907b5c588d5409c30e92e5b685c0dc28` است.

\end{document}
